\documentclass[11pt]{article}
\usepackage{amsmath,amssymb}
\usepackage[hidelinks]{hyperref}
\title{Which crew needs the recovery report?}
\date{}
\begin{document}
\maketitle

Two crews must each complete one transfer by the next physical cut.
Either local recovery command can be individually authorized and
successful in isolation. But a staging port can complete only one
transfer at that cut, and an interlock lets at most one claim succeed.
One truthful bit can restore service, but only if it reaches the crew
whose unresolved distinction changes the joint choice. Even when every
physical history has an offline assignment, no pair of local handlers
need realize all those assignments.

\section{Source, safety and service}
These are civil, one-shot recovery models, with two crews \(L,R\) and
two independent exclusive ports labelled 0 and 1. An older fallback
at each crew remains physically active through the attempted transfer.
Before the final dispatch each crew knows its own observed label,
including which ports are both lawful and physically reachable. A
source history gives a pair of labels. It contains a late loss or
displacement after preparation; controllers may not read the history
ID or the other crew's label and cannot communicate before committing
unless an explicit report is added. The admissible source support
is fixed before evaluating policies. Exactly one actual final effect
at each crew by the cut is the positive service requirement. The
interlock can safely decline one joint request, but a decline cannot
be counted toward service. There is no retry that finishes in the
last slot.

In the archive fixture, the old durable record remains official
until its one valid new write is physically recorded. In the
cooling proxy, the old fallback maintains two units at each load.
Each successful new transfer combines one local base unit and one
unit from its selected branch; a failed transfer leaves the old
supply in place within the declared window. The two proxies have
different effect checks, but share a purposely identical final
port-conflict law. They are not an independent transfer test. The
mechanisms and time scale are laboratory abstractions, not advice
for running real machinery.

\section{Exact two-port selector}
Represent each possible local observation by a vertex on its crew's
side of a bipartite graph; each admitted joint history adds an edge.
An endpoint is pinned if it has just one locally legal port. On each
edge the two selected ports must differ. Complementing R's port
number gives an equivalent condition: the normalized choice bit must
be constant along every edge. A pinned vertex now prescribes its
component's bit. Thus a lawful sure-service policy exists iff no
component contains oppositely pinned vertices (assuming no active
vertex has an empty action set). Construct one by assigning each
component its pinned value or either value if unpinned, then undoing
the R complement. An opposite-pin path is a certificate of
impossibility. The graph procedure is linear in the given observation
support; it does not by itself construct that support from arbitrary
mechanics or forecast a later unmodeled fault.

There are families with arbitrarily long necessary witnesses even
though each edge is feasible alone. Put opposite pins at the ends of
a path of length \(d\geq2\) and leave all internal vertices unpinned.
Deleting any one edge separates the pins, so no proper subset of
the \(d\) histories witnesses impossibility. The number of possible
labels grows in this family. This is a scoped source-history
certificate lower bound, not a lower bound on the memory of an
arbitrary recurring corridor controller.

\section{A report that reaches just one crew}
The path support can be modeled by a displaced cam on one of \(d\)
detents \(e=0,\ldots,d-1\). Two staggered scales give the two crews
\(\lfloor(e+1)/2\rfloor\) and \(\lfloor e/2\rfloor\); route access at
the two end labels gives the stated endpoint pins. All internal
labels leave both ports available. The model accepts these sensor
and access laws as physical premises. A separate threshold pickup
can report whether \(e\geq j\), \(0<j<d\), and may be wired to one
crew only. It produces one bit before dispatch, not a computed joint
action or an inferred exact detent. Routing a bit to a crew splits
only that crew's observation vertices according to the reported value.
In this path, the two sides of the threshold transition share just
internal vertex \(j\). The report restores sure service iff it
reaches the crew owning that vertex (or both crews) before dispatch.
If only the other crew gets it, no vertex relevant to the obstruction
splits. A late report cannot enter the choice. Both negative controls
retain the same underlying physical fault and resources.

This is a constructive test of the information interface. At a fixed
threshold the slot and scope effects are identical across report
arms. An unconsumed copy of truthful evidence is not extra authority
or usable service. An exact selector that receives the correct bit
ties the graph rule. The gain over a late or wrongly routed report
is a real modeled interface change, not a win over that matched
selector.

\section{Sharp robust service probabilities}
Now allow random choices but retain the physical interlock. The
source adversary selects the edge before random draws and may choose
the worst one for the known policy. It cannot observe the realized
seed and then choose a failing history. Safety is unconditional
inside the modeled fallback window. The probabilities below are
for completing \emph{both} effects, not for a collision occurring.

If the graph has mixed pins, let \(d\) be the length of its shortest
opposite-pin path. With a common random seed the optimal worst-history
service probability is \(1-1/d\). On a shortest path every deterministic
policy breaks at least one equality; averaging over that path gives
a history with failure probability at least \(1/d\).
For attainability take graph distances from all zero pins. Choose a
uniform threshold among \(0,\ldots,d-1\); set a vertex's normalized
bit to zero iff its distance is at most that threshold. One pins
are all at distance at least \(d\); disconnected vertices can be
regarded as at infinite distance. An edge connects adjacent or equal
distance layers, so it is broken by at most one threshold. If no
component has mixed pins, the original deterministic policy gives
probability one. A path of length one has a physically impossible
edge and probability zero regardless of information or randomness.

With independent private randomization, the value for finite
\(d\geq2\) is exactly \(1/2\). Along an opposite-pin path, the
probabilities of outputting normalized one go from zero to one;
some successive pair \(p,q\) has \(p\leq1/2\leq q\).
The vertices are owned by different crews, so their disagreement
probability is \(p(1-q)+(1-p)q\geq1/2\). Independent fair choices at
unpinned vertices attain the bound: both pinned endpoints of an edge
agree unless \(d=1\), and any edge incident to a fair unpinned vertex
succeeds with probability one half. This comparison excludes another
source of shared random state. Draws at the same crew's different
unrealized observations may be arbitrarily dependent; the argument
uses only independence between actual crew choices at each edge.

\paragraph{How much common randomness?}
For a single \(d\)-edge path, an optimal distribution over deterministic
table pairs needs at least \(d\) distinct shared outcomes. At the optimum
each edge must fail with probability \(1/d\), and a realized pair must
break exactly one edge on that path. With additional independent
private coins conditional on the shared outcome, exactly
\(\lceil d/2\rceil\) shared outcomes suffice and are necessary.
A fixed shared outcome can randomize the broken edge between two
successive positions by letting the owner of their common vertex flip
a private coin. It cannot give positive mass to breaks separated by
at least two edges. In such a case two successive interior vertices
of opposite owners change value across the two candidate breaks.
Conditional independence would permit their two tables to be combined
in the order that puts one before zero along the path. The endpoints
still require zero then one, so that realization breaks more than one
edge; it cannot occur at the optimum. Thus a shared outcome covers
at most two adjacent breaks. Group consecutive pairs (and one
singleton if needed), choosing a group with probability its size divided
by \(d\) and a fair private coin for its pair. The sharp common-support
upper bound here is only for a single simple path. Arbitrary graph
distance layers can include vertices controlled by both parties.
No assertion is made about repeated independent failures on a reused
seed, or entropy under a prescribed fair-bit generator.

\section{Comparison and limits}
\href{https://cs.uwaterloo.ca/~watrous/Papers/NonlocalGames.pdf}{Cleve,
H{\o}yer, Toner and Watrous} provide the usual two-party XOR game
framework. This model is a worst-fixed-history classical specialization
with action pins and a safety-preserving physical interlock; it makes
no claim about quantum advantage or priority of XOR formulations.
\href{https://www.sciencedirect.com/science/article/pii/S0005109815003490}{Yin
and Lafortune} are pertinent when delayed fault observations refine
a control interface. The simple exact graph solver here relies on two
ports and the observation support being supplied. A finite policy
enumerator and small same-information mixed-policy LP agree with
the stated bounds. Lean checks the Boolean normalization and two
arithmetic steps, not the provenance of the source table.

The test exercises a local-handler composition failure together with
a lawful positive effect in two source proxies. It leaves independent
mission transfer, arbitrary source-derived graph generation, a matched
advantage over a full recovery stack and inventory-wide acceptance
open. An ISS or other hardware trace with only nominal telemetry
would not by itself establish the admitted joint hidden histories.
\end{document}
